% Subsection: Finite Operations  (notes stub; content authored later)
\subsection{Finite Operations}

\begin{remark*}[Exposition]
Finite operations combine only finitely many selected subsets at a time.
They are the first algebraic layer because they mirror ordinary algebraic
manipulation: take two sets, then three, then any finite list, and ask how the
family behaves under those finite combinations.
\end{remark*}

\begin{definition}[Finite Operation on a Family]\label{def:finite-operation-on-family}
Let \(X\) be a set and let \(\mathcal{F}\subseteq\mathcal{P}(X)\).  A
\emph{finite operation on \(\mathcal{F}\)} is an operation formed from finitely
many inputs \(A_1,\dots,A_n\in\mathcal{F}\), where \(n\geq 1\), using a fixed
finite set-theoretic rule.
\end{definition}

\begin{remark*}[Standard quantified statement]
A typical finite union operation on \(\mathcal{F}\) has the form
\[
  A_1,\dots,A_n\in\mathcal{F}
  \quad\Longmapsto\quad
  \bigcup_{k=1}^{n}A_k\subseteq X,
  \qquad n\geq 1.
\]
A typical finite intersection operation has the form
\[
  A_1,\dots,A_n\in\mathcal{F}
  \quad\Longmapsto\quad
  \bigcap_{k=1}^{n}A_k\subseteq X,
  \qquad n\geq 1.
\]
\end{remark*}

\begin{remark*}[Interpretation]
Finite operations are insensitive to infinite limiting behavior.  They ask
what can be built from a bounded amount of set manipulation.  Rings and
algebras of sets will be organized around this finite level.
\end{remark*}

\begin{remark*}[Examples]
If \(A,B,C\in\mathcal{F}\), then
\[
  (A\cup B)\cap C,
  \qquad
  A\setminus(B\cap C),
  \qquad
  (A^c\cup B)\cap C
\]
are all outputs of finite operations on members of \(\mathcal{F}\).
\end{remark*}

\begin{dependencies}
\begin{itemize}
  \item \hyperref[def:family-of-subsets]{Family of Subsets}
  \item \hyperref[def:finite-infinite]{Finite and Infinite Sets}
  \item \hyperref[def:union]{Union}
  \item \hyperref[def:intersection]{Intersection}
\end{itemize}
\end{dependencies}
